\documentclass[12pt]{article}
\usepackage[left=0.75in,right=0.75in,top=0.75in,bottom=0.55in]{geometry}
\usepackage{fontspec}
\setmainfont{Times New Roman}
\usepackage[hidelinks]{hyperref}
\usepackage{parskip}
\usepackage{setspace}
\setlength{\parskip}{0.35em}
\pagestyle{empty}

\begin{document}
\onehalfspacing

Dear Waymo Hiring Team,

I am applying for the Senior Security Engineer position at Waymo. I bring more than ten years of software engineering and cybersecurity experience across critical-infrastructure and public-sector environments. My work combines software development with threat modeling, vulnerability assessment, penetration testing, secure coding remediation, malware and APT analysis, and incident response. I am drawn to roles where technical security decisions directly affect safety, reliability, and public trust---a combination central to Waymo's mission and closely aligned with the work that has shaped my career.

I began my career at KEPCO E\&C in the Nuclear Steam Supply System Development Division, developing reactor design software and engineering databases while performing vulnerability assessments and secure coding remediation for nuclear-plant systems. This work gave me a practical understanding of how software, engineering processes, and operational constraints interact where reliability is essential. It taught me to approach security as a lifecycle responsibility: understand how a system is designed and used, identify realistic attack paths, and work with technical stakeholders on effective, operationally practical mitigations.

At the Ministry of Science and ICT, I expanded that foundation across national public-sector and critical-infrastructure environments. I delivered penetration testing, red teaming, vulnerability analysis, and remediation consulting to more than 200 public institutions, including nuclear, thermal-power, railway-control, and power-system environments. In my current threat-management work, I support government organizations through malware analysis, attack-evidence collection, APT tracking, technical reporting, and remediation guidance, contributing more than 500 public cyber-threat intelligence cases each year. These responsibilities demand disciplined analysis and clear communication because findings must become concrete actions that engineering teams and decision-makers can evaluate, prioritize, and sustain.

In 2023, I initiated a research project with the Korea Internet \& Security Agency on cyber incident response for critical infrastructure. We analyzed hacking infrastructure used by international threat groups and explored how AI and data analytics could improve attack prediction, detection, and mitigation. The project reflects how I work: investigate carefully, challenge assumptions, connect evidence across systems, and communicate conclusions that support timely risk reduction.

Waymo's need for rigorous threat and risk assessment throughout the product lifecycle strongly matches this background. I can contribute experience examining complex environments, identifying threats and vulnerabilities, developing practical remediation guidance, and collaborating across engineering and operational stakeholders. I recognize that automotive cybersecurity brings specialized standards and safety interactions; I would bring a strong foundation in critical-infrastructure security and high-consequence risk analysis while deepening my expertise in ISO/SAE 21434 and autonomous-driving security. I would welcome the opportunity to help Waymo strengthen security requirements and protect the Waymo Driver throughout development and operation.

Thank you for your consideration.

Sincerely,\\
Kevin Lee

\end{document}
